\section{INTRODUCTION}
\label{sec:introduction}

Federated Learning (FL) prevents raw data sharing among clients, but joint model training can still leak information through repeated gradient exchange~\cite{zhu2019deep, so2023securing}. One-shot FL reduces communication overhead and privacy exposure by limiting communication to a single round, where intermediate gradient exchanges are not performed. However, if the transferred distillation targets are visible to the server in plaintext, the privacy risk remains~\cite{jagielski2023students}.

Differential Privacy (DP) and secure aggregation partially address this risk, with their own limitations. Strong DP noise degrades distillation quality~\cite{shao2023selective}, while secure aggregation protects individual values without preventing the server from inspecting the aggregate~\cite{kerkouche2023client}. Homomorphic encryption (HE) provides end-to-end cryptographic protection; yet direct iterative HE training of deep models is prohibitively expensive due to multiplicative depth constraints~\cite{zhang2020batchcrypt, jin2023fedml, kanpak2024cure}. This motivates a distillation-based design that is both one-shot and HE-compatible. The main technical challenge is that HE computation supports only addition and multiplication over ciphertexts, whereas standard neural network operations, e.g., ReLU, batch normalisation, softmax, are non-polynomial. Recent work addresses the stability of polynomial activations for HE-compatible networks by deriving tighter ReLU approximations~\cite{agamennone2025polynomial} and proposing boundary-loss and gradient-clipping strategies for deep polynomial training~\cite{alhossain2025training}. However, these methods target centralised inference or single-model training. Stable and efficient training of deep polynomial student models under encrypted federated supervision, where heterogeneous clients produce targets with incompatible magnitude distributions, remains an open problem.

% Also, while existing work on one-shot HE distillation focuses on encrypted inference~\cite{agamennone2025polynomial, alhossain2025training}, training deep student models efficiently under encryption remains an open problem.


For this reason, we focus on one-shot federated knowledge distillation under homomorphic encryption. In knowledge distillation, a compact \emph{student} model learns by imitating the internal representations of a stronger pre-trained \emph{teacher}, rather than training directly on labelled data. In our setting, each client acts as a teacher by training a local model on its private data and contributing encrypted supervision signals derived from that model in a single communication round. The server then uses these signals to train a privacy-preserving student entirely on ciphertexts, without ever accessing any plaintext data, weights, or activations.

% For this reason, we focus on one-shot federated knowledge distillation under homomorphic encryption: clients contribute encrypted teacher supervision in a single round, and the server trains a privacy-preserving student model entirely on ciphertexts, without ever accessing any plaintext data, weights, or activations.

\subsection{Motivation}
\label{sec:intro_motivation}

One-shot federated learning~\cite{guha2019oneshot, wang2025towards} limits client--server communication to a single round, eliminating the iterative exposure that facilitates multi-round attacks. Knowledge distillation suits this setting naturally: clients train local teachers and transfer knowledge to a server-side student in a single step. However, data heterogeneity causes inconsistent teacher logits, and as shown in Section~\ref{sec:bg_server_inference}, even this single round of plaintext transfer remains exploitable by the server. Encrypting the transferred knowledge is therefore necessary, not optional.

Homomorphic encryption (HE)~\cite{cheon2017ckks} provides end-to-end cryptographic protection without injecting noise, avoiding the utility--privacy tradeoff of DP. However, applying HE to standard iterative training is structurally expensive. The systems reviewed in Section~\ref{sec:bg_he_fl} (POSEIDON, BatchCrypt, FedSHE) all encrypt per-round gradient updates and therefore exhaust the multiplicative-depth budget of CKKS~\cite{cheon2017ckks} (further discussed in Section~\ref{sec:ckks_prelim}) in every round; bootstrapping~\cite{cheon2018bootstrapping} refreshes the budget but at substantial latency. One-shot settings resolve this tension: ciphertexts are used once, and block-independent training with fresh ciphertexts confines the depth budget to a single block of the student rather than the full network. This principle was recently validated for encrypted MLP training~\cite{pirillo2025reboot}. HE-IFD extends it to federated one-shot distillation over heterogeneous clients.

DP partially addresses plaintext leakage but degrades the distillation signal. Applying local DP causes severe accuracy drops (up to $70\%$ on CIFAR-100 at $\epsilon = 0.1$~\cite{gad2024communication}), and one-shot settings consume the entire privacy budget immediately, precluding amplification across rounds. At meaningful privacy levels ($\epsilon \leq 1$), DP reduces FedMD to near random-guessing accuracy~\cite{sun2021fedmdnfdp}. HE avoids this tradeoff entirely.

\subsection{Our Approach}
We present \textbf{HE-IFD} (\textbf{H}omomorphic \textbf{E}ncryption-\textbf{I}ntermediate \textbf{F}eature \textbf{D}istillation), a framework in which individual client teacher outputs are protected during knowledge transfer and all subsequent server-side computation. The protocol operates in three phases:

\begin{enumerate}
\item \textbf{Local training.} Each client trains a standard teacher model (e.g., ResNet-18) on their private data partition, starting from a shared random initialisation.
\item \textbf{Encrypted upload.} Each client runs its teacher model on its own data and records the inputs and outputs at each layer group boundary. We refer to these groups as \emph{blocks}. These input--output pairs tell the server what each block of the teacher produces, without revealing the raw data. Before encrypting, clients jointly agree on per-channel normalisation statistics so that features from different clients are on a comparable scale. The normalised pairs are then encrypted under Cheon--Kim--Kim--Song (CKKS) scheme and uploaded to the server in a single round; no further client interaction is needed.
\item \textbf{Server-side encrypted distillation.} The server trains a HE-compatible architecture in which non-polynomial operations such as ReLU and batch normalisation are replaced with polynomial equivalents that can be evaluated directly on ciphertexts (e.g., PolyResNet-18) for the student model on the server, block by block, using only encrypted inputs and targets. Student weights, activations, gradients, and loss values are all ciphertext throughout training.
\end{enumerate}

The server never observes any individual teacher output, student weight, or intermediate activation in plaintext. It operates as a blind compute delegate that executes the distillation protocol on encrypted data. Four properties of the protocol class determine the design and govern what can be proven about it; each is the structural answer to a distinct concern raised by the rejected version's review cycle.

\par\noindent\textbf{(C1) HE depth budget for an end-to-end encrypted student update.} A naive encrypted SGD step that forward-propagates through encrypted weights accumulates one multiplicative level per affine layer plus the degree of every polynomial activation, which exceeds the level chain available in any practical CKKS parameterisation long before the student converges. We avoid the naive construction. The encrypted state at every step is a \emph{linear accumulator} over teacher-induced gradient contributions, $\langle\theta_E\rangle = \langle\theta_0^\star\rangle + \sum_t \eta\,\langle g_t\rangle$, where each $\langle g_t\rangle$ is computed from the encrypted teacher signal applied to a plaintext student state. The per-step depth is at most three multiplicative levels (residual carry-over, scalar-by-ciphertext for the learning rate, addition to the accumulator), and the dominant ciphertext-by-ciphertext cost is confined to the depth-3 ensemble target construction described in Section~\ref{sec:phase2}; the full protocol fits inside a single CKKS level chain at $\log N \in \{14, 15\}$ with no bootstrapping. The structural justification for the plaintext student state during training, and for what is and is not allowed to be threshold-decrypted, is the binding invariant developed in Section~\ref{sec:threat_binding}.

\par\noindent\textbf{(C2) $\beta/\lambda$ ensemble boost without division under HE.} Pooling per-client encrypted teacher logits into an ensemble target with confidence-aware reweighting requires both a per-client confidence weight $\alpha_i^\beta$ and a per-row data-boost factor $1+\lambda V_k$ on the encrypted variance across teachers. Both knobs are constructed in CKKS without any encrypted division: the $\beta$-power is taken on the plaintext scalar $\alpha_i$ and applied as a scalar-by-ciphertext multiplication, and the per-row variance is built as $\langle V_k\rangle = \sum_i w_i \langle T_{i,k}\rangle^2 - \langle\bar T_k\rangle^2$ at depth 1. The resulting ensemble target is at most three multiplicative levels deep regardless of $N$. The construction is detailed in Section~\ref{sec:phase2}; the same section formalises why the depth bound is independent of $N$ and why a per-class coverage-aware variant of $\beta$ is the natural extension at the lowest-$\alpha$ heterogeneity settings.

\par\noindent\textbf{(C3) Binding invariant under $N{-}1$ collusion.} The cryptographic primitive is the multiparty variant of CKKS in the threshold instantiation at $t{=}N$: all $N$ secret-key shares must concur to recover any plaintext. The choice $t{=}N$ rather than $t < N$ is deliberate, because any $t < N$ admits a coalition of $t{-}1$ semi-honest clients who can decrypt the per-client logit ciphertexts and the aggregated ensemble target, which under the present threat model directly enables subtraction attacks against the remaining honest client. The protocol's central security property, that the only ciphertext ever subjected to threshold decryption is the final trained student $\theta_E$, is the binding invariant, formalised in Section~\ref{sec:threat_binding}. With the invariant, the adversary's view (server in coalition with up to $N{-}1$ clients) consists of IND-CPA ciphertexts, plaintext metadata (probe set, ciphertext sizes, timing), and the released student alone.

\par\noindent\textbf{(C4) Post-release statistical-query-floor mitigation.} The binding invariant prevents direct decryption attacks on intermediate state, but it does not rule out a derived form, utility-coercion-into-privacy-amplification: if $N{-}1$ colluding clients upload all-zero logit ciphertexts, the encrypted ensemble target collapses to a multiple of the single honest client's logits and the released student becomes a single-teacher distillation of that honest teacher. The attack does not breach the invariant; it inflates the unavoidable statistical-query floor on $\theta_E$. The protocol admits two independent defensive knobs: each client trains its teacher with DP-SGD at a per-client budget $(\varepsilon_T, \delta_T)$ that survives by post-processing into the released student, and each client may add per-row Gaussian noise to its logit vector before encryption, which surfaces at full magnitude under the all-zeros amplification and averages to $\sigma_P/\sqrt N$ in the honest case. The composition is bounded by the standard Gaussian-DP RDP accountant. The construction is developed in Section~\ref{sec:threat_binding}.

Together, these four properties (the depth budget for the linear accumulator, the division-free $\beta/\lambda$ ensemble boost, the binding invariant on threshold decryption, and the SQ-floor mitigation by DP-SGD teachers plus per-row Gaussian noise) determine the design of the protocol. Section~\ref{sec:methodology} develops each of them; Section~\ref{sec:experiments} reports the resulting accuracy on MNIST and CIFAR-10 against the strongest published one-shot baselines, and shows that the encrypted student consistently exceeds the mean individual teacher accuracy across heterogeneity settings, with the largest margins at the smallest $\alpha$.


% A key technical difficulty is training an HE-compatible student, which uses polynomial activations and HE-compatible normalisation layers instead of ReLU and batch normalisation, from encrypted supervision. Standard end-to-end distillation diverges for this architecture. We solve this with training the model sequentially, one block at a time, where each block is an independent sub-network spanning a few layers, with each block receiving fresh ciphertexts, combined with magnitude-regularised normalised MSE loss and student-input training. This approach yields stable convergence and competitive accuracy across our CIFAR-10 and Fashion-MNIST experiments.


\subsection{Contributions}
Our contributions are as follows:

\begin{enumerate}
\item We propose, for the first time, progressive stacking for training HE-compatible student models fully under encryption: a block-by-block procedure with magnitude-regularised normalised MSE loss that achieves stable distillation of a full PolyResNet-18 student (11.17M parameters) from encrypted teacher features, where standard end-to-end distillation diverges.

\item We present a novel, fully encrypted one-shot distillation protocol in which clients upload encrypted per-block feature pairs in a single communication round and the server performs all distillation computation (forward pass, loss evaluation, backward pass, and weight update) entirely on ciphertexts, without exposing any plaintext data at any point during training.

\item We introduce client-side collaborative normalisation, a lightweight technique in which each client normalises its teacher outputs to zero mean and unit variance per channel before encryption, which eliminates magnitude incompatibility across heterogeneous teachers and prevents training divergence for $N \geq 4$ clients.

\item We provide a comprehensive empirical evaluation of HE-IFD across three architectures (ResNet-18, SimpleCNN, ViT-B/32), client counts $N \in \{1, 2, 4, 8, 16, 32\}$, and heterogeneity levels $\alpha \in \{0.1, 1.0\}$, including accuracy comparisons against differential privacy baselines, membership inference attack analysis, and communication and computation cost analysis.
\end{enumerate}

%The remainder of this paper is organized as follows. Section~\ref{sec:background} surveys federated learning, knowledge distillation, privacy mechanisms, and homomorphic encryption. Section~\ref{sec:methodology} presents the HE-IFD framework. Section~\ref{sec:experiments} reports experimental results. Section~\ref{sec:conclusion} concludes.

    

% \section{INTRODUCTION}
% \label{sec:introduction}

% Federated Learning (FL) prevents raw data sharing among clients, but joint model training can still leak information through repeated gradient exchange~\cite{zhu2019deep, so2023securing}. One-shot FL reduces communication overhead and privacy exposure by limiting communication to a single round, where intermediate gradient exchanges are not performed. However, if the transferred distillation targets are visible to the server in plaintext, the privacy risk remains~\cite{jagielski2023students}.

% Differential Privacy (DP) and secure aggregation partially address this risk, with their own limitations. Strong DP noise degrades distillation quality~\cite{shao2023selective}, while secure aggregation protects individual values without preventing the server from inspecting the aggregate~\cite{kerkouche2023client}. Homomorphic encryption (HE) provides end-to-end cryptographic protection; yet direct iterative HE training of deep models is prohibitively expensive due to multiplicative depth constraints~\cite{zhang2020batchcrypt, jin2023fedml, kanpak2024cure}. This motivates a distillation-based design that is both one-shot and HE-compatible. The main technical challenge is that HE computation supports only addition and multiplication over ciphertexts, whereas standard neural network operations, e.g., ReLU, batch normalisation, softmax, are non-polynomial. Recent work addresses the stability of polynomial activations for HE-compatible networks---deriving tighter ReLU approximations~\cite{agamennone2025polynomial} and proposing boundary-loss and gradient-clipping strategies for deep polynomial training~\cite{alhossain2025training}, but these methods target centralised inference or single-model training. Stable and efficient training of deep polynomial student models under encrypted federated supervision, where heterogeneous clients produce targets with incompatible magnitude distributions, remains an open problem.

% % Also, while existing work on one-shot HE distillation focuses on encrypted inference~\cite{agamennone2025polynomial, alhossain2025training}, training deep student models efficiently under encryption remains an open problem.




% For this reason, we focus on one-shot federated knowledge distillation under homomorphic encryption. In knowledge distillation, a compact \emph{student} model learns by imitating the internal representations of a stronger pre-trained \emph{teacher}, rather than training directly on labelled data. In our setting, each client acts as a teacher by training a local model on its private data and contributing encrypted supervision signals derived from that model in a single communication round. The server then uses these signals to train a privacy-preserving student entirely on ciphertexts, without ever accessing any plaintext data, weights, or activations.

% % For this reason, we focus on one-shot federated knowledge distillation under homomorphic encryption: clients contribute encrypted teacher supervision in a single round, and the server trains a privacy-preserving student model entirely on ciphertexts, without ever accessing any plaintext data, weights, or activations.

% \subsection{Our Approach}
% We present \textbf{HE-IFD} (\textbf{H}omomorphic \textbf{E}ncryption-\textbf{I}ntermediate \textbf{F}eature \textbf{D}istillation), a framework in which individual client teacher outputs are protected during knowledge transfer and all subsequent server-side computation. The protocol operates in three phases:

% \begin{enumerate}
% \item \textbf{Local training.} Each client trains a standard teacher model (e.g., ResNet-18) on their private data partition, starting from a shared random initialisation.
% \item \textbf{Encrypted upload.} Each client runs its teacher model on its own data and records the inputs and outputs at each layer group boundary — we call these groups \emph{blocks}. These input--output pairs tell the server what each block of the teacher produces, without revealing the raw data. Before encrypting, clients jointly agree on per-channel normalisation statistics so that features from different clients are on a comparable scale. The normalised pairs are then encrypted under CKKS and uploaded to the server in a single round; no further client interaction is needed.
% \item \textbf{Server-side encrypted distillation.} The server trains a HE-compatible architecture in which non-polynomial operations such as ReLU and batch normalisation are replaced with polynomial equivalents that can be evaluated directly on ciphertexts (e.g., PolyResNet-18) in the student model in server, block by block, using only encrypted inputs and targets. Student weights, activations, gradients, and loss values are all ciphertext throughout training.
% \end{enumerate}

% The server never observes any individual teacher output, student weight, or intermediate activation in plaintext. It operates as a blind compute delegate that executes the distillation protocol on encrypted data.

% A key technical difficulty is training an HE-compatible student — which replaces ReLU with learnable degree-2 polynomial activations and batch normalisation with static per-channel affine layers — from encrypted intermediate supervision. Several compounding challenges make standard distillation approaches fail.

% \par\noindent\textbf{Polynomial magnitude explosion.} The quadratic activation $ax^2 + bx$ amplifies feature magnitudes at every composition. An input of magnitude 5 produces roughly 11 after one block, 44 after two; across five blocks this can reach $10^4$, causing gradient overflow. The problem is exacerbated in the multi-client setting, where teachers trained on heterogeneous data partitions produce features with incompatible magnitude ranges, and the HE-compatible replacement for batch normalisation — a static per-channel scale-and-shift — cannot dynamically compensate.

% \par\noindent\textbf{Training--distillation distribution gap.} Distilling block by block avoids CKKS depth accumulation, but training block $k$ on teacher-produced inputs creates a mismatch: at inference, block $k$ receives student-produced inputs from the frozen prefix. Without correction this collapses the composed model toward random accuracy. We address this through trainable per-channel affine bridges, initialised from uploaded feature statistics, followed by server-side sequential refinement that exposes each block to its actual inference-time distribution — without any additional client communication.

% \par\noindent\textbf{Scale-anchored distillation loss.} Standard MSE is sensitive to the sign asymmetry between polynomial activations and ReLU features and allows unbounded magnitude drift that amplifies through frozen downstream blocks. Our loss combines channel-normalised MSE, which aligns feature \textit{shape} independently of scale, with a log-ratio penalty $(\log,\sigma_{\hat{f}}/\sigma_{\tilde{f}})^2$ that anchors the student's feature \textit{magnitude} to the teacher's. Unlike a plain $\ell_2$ scale penalty, the log-ratio form penalises multiplicative over- and under-scaling symmetrically and is scale-free.

% Together these mechanisms yield stable convergence and competitive accuracy across our CIFAR-10 and Fashion-MNIST experiments under non-IID federated partitioning.

% % A key technical difficulty is training an HE-compatible student, which uses polynomial activations and HE-compatible normalisation layers instead of ReLU and batch normalisation, from encrypted supervision. Standard end-to-end distillation diverges for this architecture. We solve this with training the model sequentially, one block at a time, where each block is an independent sub-network spanning a few layers, with each block receiving fresh ciphertexts, combined with magnitude-regularised normalised MSE loss and student-input training. This approach yields stable convergence and competitive accuracy across our CIFAR-10 and Fashion-MNIST experiments.


% \subsection{Contributions}
% Our contributions are as follows:

% \begin{enumerate}
% \item We propose, for the first time, progressive stacking for training HE-compatible student models fully under encryption: a block-by-block procedure with magnitude-regularised normalised MSE loss that achieves stable distillation of a full PolyResNet-18 student (11.17M parameters) from encrypted teacher features, where standard end-to-end distillation diverges.

% \item We present a novel fully encrypted one-shot distillation protocol in which clients upload encrypted per-block feature pairs in a single communication round and the server performs all distillation computation (forward pass, loss evaluation, backward pass, and weight update) entirely on ciphertexts, without exposing any plaintext data at any point during training.

% \item We introduce client-side collaborative normalisation, a lightweight technique in which each client normalises its teacher outputs to zero mean and unit variance per channel before encryption, which eliminates magnitude incompatibility across heterogeneous teachers and prevents training divergence for $N \geq 4$ clients.

% \item We provide a comprehensive empirical evaluation of HE-IFD across client counts ($N \in \{1, 2, 4, 8, 16\}$) and heterogeneity levels ($\alpha \in \{0.1, 0.5, 1.0, 100.0\}$), reporting accuracy, a privacy analysis against membership inference attacks, and a comparison with differential privacy baselines.
% \end{enumerate}

% %The remainder of this paper is organized as follows. Section~\ref{sec:background} surveys federated learning, knowledge distillation, privacy mechanisms, and homomorphic encryption. Section~\ref{sec:methodology} presents the HE-IFD framework. Section~\ref{sec:experiments} reports experimental results. Section~\ref{sec:conclusion} concludes.

